% GENERATED FILE. Edit canonical lessons, then run npm run generate:books.
% canonical-source-hash: fnv1a64:61a01e2acc342d1e
% canonical-lessons: PT-C340-duvidar, PT-C340-cetico, PT-C340-negar, PT-C340-absurdo, PT-C340-discutivel

\chapter{Tenho as minhas dúvidas}
\label{ch:pt-tenho-as-minhas-duvidas}
\hlchaptermodality{340}

\begin{chapteropening}
\textbf{By the end of this chapter:} \emph{I can express doubt, deny something and call an idea absurd or debatable.}

Complete the last lesson of chapter 340: I can express doubt, deny something and call an idea absurd or debatable.
\end{chapteropening}

\section[duvidar]{duvidar — to doubt}

\label{lesson:PT-C340-duvidar}



\begin{quote}
Before the new one: say the Portuguese for the consensus, then the Portuguese for to give in.
\end{quote}

\subsection*{You'll want to know: duvidar}
\textbf{duvidar} — \textquotedblleft{}to doubt\textquotedblright{}.

\begin{quote}
Duvido que ele chegue a horas.
\end{quote}

\begin{quote}
\emph{I doubt he will arrive on time.}
\end{quote}

\emph{Duvidar que} takes the subjunctive: \emph{duvido que venha}. \emph{Duvidar de alguém} — to doubt someone.

\begin{cousinweb}
\emph{duvidar} is inherited from Latin \emph{dubitāre} \textquotedblleft{}to doubt, to waver\textquotedblright{} (Lewis \& Short), through Old Portuguese \emph{dovidar} (EtymDB).
\end{cousinweb}

\begin{grammarlens}[title={What you've built}]
One more word for telling, explaining and describing at B1.
\end{grammarlens}

\subsection*{Your turn}
\begin{itemize}
\raggedright
  \item \emph{Say it:} \emph{duvidar}
  \item \emph{Say it:} \emph{duvidar}, once more
  \item \emph{Say it:} a sentence of your own with \emph{duvidar}
  \item \emph{Recall:} say the Portuguese for the consensus, then the Portuguese for to give in, then say \emph{duvidar} again
\end{itemize}

\begin{tcolorbox}[breakable,colback=teal!4,colframe=teal!35!black,title={Before you move on}]
What does duvidar mean? (\textquotedblleft{}To doubt\textquotedblright{}.) Say it once more.
\end{tcolorbox}
\section[cético]{cético — sceptical}

\label{lesson:PT-C340-cetico}



\begin{quote}
Before the new one: say the Portuguese for to give in, then the Portuguese for to doubt.
\end{quote}

\subsection*{You'll want to know: cético}
\textbf{cético} — \textquotedblleft{}sceptical\textquotedblright{}.

\begin{quote}
Estou cético: essa promessa parece boa demais.
\end{quote}

\begin{quote}
\emph{I am sceptical: that promise seems too good.}
\end{quote}

Many speakers in Portugal pronounce the \emph{p} and write \emph{céptico}; the agreement accepts both spellings. Feminine \emph{cética}.

\begin{cousinweb}
\emph{cético} is a learned word that goes back to Greek \emph{skeptikós} \textquotedblleft{}inquiring, reflective\textquotedblright{}, built on \emph{sképtesthai} \textquotedblleft{}to consider, to examine\textquotedblright{}; Latin had \emph{Sceptici} for the philosophers of that school (Lewis \& Short).
\end{cousinweb}

\begin{grammarlens}[title={What you've built}]
One more word for telling, explaining and describing at B1.
\end{grammarlens}

\subsection*{Your turn}
\begin{itemize}
\raggedright
  \item \emph{Say it:} \emph{cético}
  \item \emph{Say it:} \emph{cético}, once more
  \item \emph{Say it:} a sentence of your own with \emph{cético}
  \item \emph{Recall:} say the Portuguese for to give in, then the Portuguese for to doubt, then say \emph{cético} again
\end{itemize}

\begin{tcolorbox}[breakable,colback=teal!4,colframe=teal!35!black,title={Before you move on}]
What does cético mean? (\textquotedblleft{}Sceptical\textquotedblright{}.) Say it once more.
\end{tcolorbox}
\section[negar]{negar — to deny; to refuse}

\label{lesson:PT-C340-negar}



\begin{quote}
Before the new one: say the Portuguese for to doubt, then the Portuguese for sceptical.
\end{quote}

\subsection*{You'll want to know: negar}
\textbf{negar} — \textquotedblleft{}to deny; to refuse\textquotedblright{}.

\begin{quote}
Ele nega tudo, mas toda a gente viu.
\end{quote}

\begin{quote}
\emph{He denies everything, but everyone saw.}
\end{quote}

\emph{Negar que} takes the subjunctive. Also to refuse: \emph{negar ajuda}.

\begin{cousinweb}
\emph{negar} is inherited from Latin \emph{negāre} \textquotedblleft{}to say no, to deny\textquotedblright{} (Lewis \& Short), through Old Portuguese \emph{negar} (EtymDB).
\end{cousinweb}

\begin{grammarlens}[title={What you've built}]
One more word for telling, explaining and describing at B1.
\end{grammarlens}

\subsection*{Your turn}
\begin{itemize}
\raggedright
  \item \emph{Say it:} \emph{negar}
  \item \emph{Say it:} \emph{negar}, once more
  \item \emph{Say it:} a sentence of your own with \emph{negar}
  \item \emph{Recall:} say the Portuguese for to doubt, then the Portuguese for sceptical, then say \emph{negar} again
\end{itemize}

\begin{tcolorbox}[breakable,colback=teal!4,colframe=teal!35!black,title={Before you move on}]
What does negar mean? (\textquotedblleft{}To deny; to refuse\textquotedblright{}.) Say it once more.
\end{tcolorbox}
\section[absurdo]{absurdo — absurd}

\label{lesson:PT-C340-absurdo}



\begin{quote}
Before the new one: say the Portuguese for sceptical, then the Portuguese for to deny.
\end{quote}

\subsection*{You'll want to know: absurdo}
\textbf{absurdo} — \textquotedblleft{}absurd\textquotedblright{}.

\begin{quote}
É absurdo pagar tanto por um café.
\end{quote}

\begin{quote}
\emph{It is absurd to pay so much for a coffee.}
\end{quote}

Also a noun: \emph{é um absurdo!} — it is outrageous!

\begin{cousinweb}
\emph{absurdo} is a learned borrowing from Latin \emph{absurdus} \textquotedblleft{}out of tune, incongruous\textquotedblright{} (EtymDB), built from \emph{ab} and \emph{surdus} \textquotedblleft{}deaf, dull-sounding\textquotedblright{} (Lewis \& Short).
\end{cousinweb}

\begin{grammarlens}[title={What you've built}]
One more word for telling, explaining and describing at B1.
\end{grammarlens}

\subsection*{Your turn}
\begin{itemize}
\raggedright
  \item \emph{Say it:} \emph{absurdo}
  \item \emph{Say it:} \emph{absurdo}, once more
  \item \emph{Say it:} a sentence of your own with \emph{absurdo}
  \item \emph{Recall:} say the Portuguese for sceptical, then the Portuguese for to deny, then say \emph{absurdo} again
\end{itemize}

\begin{tcolorbox}[breakable,colback=teal!4,colframe=teal!35!black,title={Before you move on}]
What does absurdo mean? (\textquotedblleft{}Absurd\textquotedblright{}.) Say it once more.
\end{tcolorbox}
\section[discutível]{discutível — debatable, questionable}

\label{lesson:PT-C340-discutivel}



\begin{quote}
Before the new one: say the Portuguese for to deny, then the Portuguese for absurd.
\end{quote}

\subsection*{You'll want to know: discutível}
\textbf{discutível} — \textquotedblleft{}debatable, questionable\textquotedblright{}.

\begin{quote}
Essa decisão foi discutível.
\end{quote}

\begin{quote}
\emph{That decision was debatable.}
\end{quote}

Built on \emph{discutir} + \emph{-ível}: something you can argue about.

\begin{cousinweb}
\emph{discutível} is built on \emph{discutir} \textquotedblleft{}to discuss\textquotedblright{} with the suffix \emph{-ível}; \emph{discutir} comes from Latin \emph{discutere} \textquotedblleft{}to strike asunder, to shatter\textquotedblright{} (Lewis \& Short), which later meant \textquotedblleft{}to examine\textquotedblright{}.
\end{cousinweb}

\begin{grammarlens}[title={What you've built}]
One more word for telling, explaining and describing at B1.
\end{grammarlens}

\subsection*{Your turn}
\begin{itemize}
\raggedright
  \item \emph{Say it:} \emph{discutível}
  \item \emph{Say it:} \emph{discutível}, once more
  \item \emph{Say it:} a sentence of your own with \emph{discutível}
  \item \emph{Recall:} say the Portuguese for to deny, then the Portuguese for absurd, then say \emph{discutível} again
\end{itemize}

\begin{tcolorbox}[breakable,colback=teal!4,colframe=teal!35!black,title={Before you move on}]
What does discutível mean? (\textquotedblleft{}Debatable, questionable\textquotedblright{}.) Say it once more.
\end{tcolorbox}
